%%%%%%%%%%%%%%%%%%%%%%%%%%%%%%%%%
%
%       EXERCÍCIO
%
%%%%%%%%%%%%%%%%%%%%%%%%%%%%%%%%%

\ifdefstring{\atividade}{prova}{%
    \renewcommand{\valorquestao}{\ValorQ\ pontos}
}%

\begin{exercicioBanco}[\valorquestao]
A tabela a seguir apresenta a distribuição de frequências relativas aos tempos de atendimento (em minutos) registrados em um posto de saúde:

\begin{center}
\begin{tabular}{|c|c|}
\hline
Tempo (min) & Frequência ($f_i$) \\ \hline
$0 \vdash 10$ & 3 \\ \hline
$10 \vdash 20$ & 5 \\ \hline
$20 \vdash 30$ & 8 \\ \hline
$30 \vdash 40$ & 4 \\ \hline
$40 \vdash 50$ & 2 \\ \hline
\end{tabular}
\end{center}

Com base nessa tabela, determine:

\begin{enumerate}[(a)]
    \item A média (\(\mu\)) do tempo de atendimento.
    \item O primeiro quartil ($Q_1$) e o terceiro quartil ($Q_3$).
    \item A mediana ($\theta$).
    \item A moda ($Mo$) calculada pela fórmula de Czuber.
\end{enumerate}
\end{exercicioBanco}
